\documentclass[aspectratio=169]{beamer}

% Compilar con XeLaTeX
\usepackage{fontspec}
\usepackage{xcolor}
\usepackage{graphicx}
\usepackage{amsmath}

% Branding Académico (Obligatorio)
\definecolor{AzulMarino}{HTML}{1A365D}
\definecolor{GrisPizarra}{HTML}{4A5568}
\definecolor{VerdeAzulado}{HTML}{319795}
\definecolor{FondoBlanco}{HTML}{F7FAFC}

% Configuracion de Beamer
\setbeamercolor{background canvas}{bg=FondoBlanco}
\setbeamercolor{title}{fg=AzulMarino}
\setbeamercolor{frametitle}{fg=AzulMarino}
\setbeamercolor{normal text}{fg=GrisPizarra}
\setbeamercolor{itemize item}{fg=VerdeAzulado}
\setbeamercolor{block title}{bg=AzulMarino, fg=white}
\setbeamercolor{block body}{bg=white, fg=GrisPizarra}
\setbeamercolor{alertblock title}{bg=VerdeAzulado, fg=white}
\setbeamercolor{alertblock body}{bg=white, fg=GrisPizarra}

% Fuentes
\setmainfont{Inter}
\setsansfont{Montserrat}
\setmonofont{JetBrains Mono}

\title{Modelos Supervisados de Clasificación Clásica}
\subtitle{k-NN y Regresión Logística Aplicados a la Ciberseguridad}
\author{Machine Learning para Ciberseguridad}
\date{Semana 3}

\begin{document}

\begin{frame}
    \titlepage
\end{frame}

\begin{frame}{Objetivo de la Sesión}
    \begin{itemize}
        \item Comprender el algoritmo \textbf{k-Nearest Neighbors (k-NN)} y sus métricas de distancia.
        \item Entender la función Sigmoide y cómo la \textbf{Regresión Logística} estima probabilidades.
        \item Aplicar estos conceptos para detectar fraudes o Phishing.
        \item \textcolor{VerdeAzulado}{\textbf{Analizar el costo de negocio entre Falsos Positivos y Falsos Negativos.}}
    \end{itemize}
\end{frame}

\begin{frame}{k-Nearest Neighbors (k-NN)}
    \begin{block}{La Analogía del Vecindario}
        "Dime con quién andas y te diré quién eres."
        Si analizamos un nuevo sitio web y sus 5 "vecinos" más parecidos estadísticamente resultan ser sitios de \textit{Phishing}, el modelo clasificará este nuevo sitio como \textit{Phishing}.
    \end{block}
    
    \textbf{Distancia Euclídea ($L_2$):}
    $$d(\mathbf{x}, \mathbf{x}_i) = \sqrt{\sum_{j=1}^{d} (x^{(j)} - x_i^{(j)})^2}$$
\end{frame}

\begin{frame}{Regresión Logística y la Función Sigmoide}
    \begin{alertblock}{La Rampa Clasificadora}
        Imagina un tobogán liso. Los atributos sospechosos (falta de SSL, IP expuesta) suman "peso". Si la evidencia es pesada, te deslizas hacia la zona de \textit{Phishing} (probabilidad cercana a 1).
    \end{alertblock}
    
    La función sigmoide mapea cualquier número real a una probabilidad entre 0 y 1:
    $$\sigma(z) = \frac{1}{1 + e^{-z}}$$
\end{frame}

\begin{frame}{Matriz de Confusión en Ciberseguridad}
    \begin{table}[]
        \centering
        \begin{tabular}{|l|l|l|}
        \hline
        \textbf{} & \textbf{Predicción: Phishing} & \textbf{Predicción: Legítimo} \\ \hline
        \textbf{Real: Phishing} & Verdadero Positivo (VP) & \textcolor{red}{Falso Negativo (FN)} \\ \hline
        \textbf{Real: Legítimo} & \textcolor{orange}{Falso Positivo (FP)} & Verdadero Negativo (VN) \\ \hline
        \end{tabular}
    \end{table}
    
    \begin{itemize}
        \item \textbf{FP:} Usuario bloqueado por error (Fricción de Negocio).
        \item \textbf{FN:} Ataque exitoso que no detectamos (Brecha de Seguridad Catastrófica).
    \end{itemize}
\end{frame}

\end{document}
